\input{technique-note.tex}
\title{Integration Bee --- Unit 07:\\ Absolute Values \& Reverse Product Rule}
\author{UK Uni Integration Bee prep\\ Official: integration.soc.srcf.net\\ Yang Topics Guide 2nd Ed}
\date{}
\begin{document}
\maketitle
\section*{Status}
Concept teaching only. Tags: \texttt{absolute-value}, \texttt{reverse-product}, \texttt{yang:1.12}, \texttt{yang:1.13}.
\section{Official anchors}
\begin{itemize}
  \item Yang §§1.12 Absolute Values, 1.13 Reverse Product Rule.
  \item Related to parts / product rule (Syllabus §1).
\end{itemize}
\section{Learning goals}
\begin{enumerate}
  \item Split absolute-value integrands at zeros.
  \item Spot integrands that are exactly $(fg)'$.
\end{enumerate}
\section{Pedagogy}
Worked $\to$ faded $\to$ interleaved. See \texttt{pedagogy/README.md}.

\section{Worked abs --- \texttt{U07-D01} (\texttt{yang:1.12})}
\[
\int_{-1}^{1}|x|\,dx = 2\int_0^1 x\,dx = 1.
\]

\section{Worked reverse product --- \texttt{U07-D04} (\texttt{yang:1.13})}
\[
\int_0^{\pi/2}(\sin x + x\cos x)\,dx = \big[x\sin x\big]_0^{\pi/2}=\frac\pi2.
\]
(Recognise derivative of $x\sin x$.)

\section{Faded --- \texttt{U07-D05}}
\[
\int_0^{\pi/2} e^x(\sin x+\cos x)\,dx.
\]
Scaffold: derivative of $e^x\sin x$.

\section{Practice}
\texttt{U07-D01}\ldots\texttt{U07-D08} in \texttt{drills.tex}; \texttt{verify/run.sh}.
\end{document}
